\chapter{Discussion}
\label{ch:discussion}

Chapter~\ref{ch:results} reported what the measurements show. This chapter asks what they mean. It sets out the
interpretation the evidence supports, answers the objection that novelty must surely matter, reconciles the
null with a literature of positive results, draws out what follows for practice, and states the bounds within
which the claim holds.

\section{What the result means}
\label{sec:disc-meaning}
Neither hypothesis survives. H1 and H2 both fail the bar set in advance in Section~\ref{sec:meth-success}, on
both hold-out schemes and for both targets (Section~\ref{sec:headline}). H0 is what the data support: transfer
is not distance-driven.

How they fail matters as much as that they fail. The one coefficient that reached significance ran in the
wrong direction, and dissolved under a size control. The follow-up difficulty pivot behaved the same way, its
apparent effect routing entirely through animal size. So the only label-free property that predicts detection
is the very quantity the analysis introduced as a nuisance to control for.

The detection--association decomposition, the intellectual core of the design, produced a second and quieter
finding. Association is a near-degenerate target on this data. Of the $346$ positive cells, $268$ ($77\%$) hold
at most one annotated animal per video, so there is only one identity to keep. Across all positive cells
\texttt{pAssA} correlates with \texttt{pDetA} at $0.97$ and is exactly equal to it in $73\%$, leaving little
identity-specific variance for any feature to predict. The decomposition is therefore a detection story, and
the association null is reported with those divergence figures as its explanation rather than as a separate
failed search.

\subsection{Novelty at the extremes versus graded prediction}
\label{sec:disc-coarse}
A natural objection is that novelty surely matters. A familiar species is detected better than a totally alien
one, so how can a novelty distance be null?

Both statements are true, because they concern different resolutions of the same axis. The coarse claim, that
the extreme tails differ, is real and never in dispute. The null reported here is the far stronger graded claim
that H1 actually tests: that a continuous distance predicts where on the detection scale an unseen species
lands, out of sample, beyond the size confound, and well enough to gate a deployment.

The two come apart because the informative variance lives in the middle of the range rather than at the poles.
The easy extremes are already predicted by the mean, with common large animals near $1$ and hopeless cases near
$0$. A coarse alien-versus-familiar signal therefore adds nothing exactly where prediction would have value. A
genuine extreme-tail effect can coexist with a null graded predictor without contradiction, and it is the
graded, held-out, size-controlled predictor that a practitioner would need.

\section{Why the null does not contradict prior positives}
\label{sec:disc-related}
The result sits against two literatures of positive findings: one that predicts model performance without
labels, and one that shows backgrounds drive camera-trap failure. Neither is contradicted here, and it is worth
saying precisely why.

\subsection{Label-free performance estimation}
Two axes separate this study from every positive result in Section~\ref{sec:bg-predictable}.

The first is the signal. Those estimators all read the model's own behaviour on the target, after running it.
This work forecasts from external distances known before a single frame is processed. That is the strictly
harder regime, and for a practitioner deciding whether to spend compute, the more useful one.

The second is the output. Every predecessor targets a scalar classifier accuracy, or the quality of a static
detector or segmenter. None concerns identity maintained over time, and none decomposes its estimate. This is
the first study to put the question to a promptable video tracker, split into detection and association.

AutoEval~\cite{autoeval} is the closest methodological cousin, and the strength of its positive sharpens the
contrast. It reports a near-perfect rank correlation between a Fr\'echet distribution-shift distance and
classifier accuracy. That distance, though, is computed from the model's own features after the target images
have passed through the network. We tested AutoEval's exact functional in our own regime --- Fr\'echet and MMD
distances on cached DINOv2 embeddings, computed before running SAM\,3 --- and it stayed null
(Section~\ref{sec:robustness}). The failure is therefore not an artefact of a lazy nearest-prototype
distance. Our nearest analogue does worse still, running in the opposite direction before dissolving under a
size control.

One bridge makes the boundary concrete. Accuracy-on-the-Line~\cite{aotl} was validated on iWildCam, the same
camera-trap shift this dissertation cites through WILDS and Terra Incognita~\cite{wilds,terra}. The exact
wildlife shift where a classifier's out-of-distribution accuracy is predictable label-free is where our
before-running distances are not.

\subsection{Background bias}
The environment axis inherits its foreground/background separation from a literature that finds backgrounds
emphatically do matter (Section~\ref{sec:bg-cameratrap}), so its null needs reconciling.

Two distinctions dissolve the tension. The first is causal reliance versus predictive distance.
\citet{xiao}, \citet{elephant} and PanAf-FGBG~\cite{panaf} all intervene on the background and watch
performance move. That is a statement about how a model uses context. Ours is the strictly weaker, correlational question of whether a
scalar scene distance forecasts an untouched model's per-cell score. A model can lean causally on context while
its score stays unpredictable from a one-number distance. The null does not deny their causal findings. It
locates a regime where the distance shortcut fails.

The second distinction is why the promptable regime should be that regime, and the same literature supplies the
mechanism. Those models must infer the category, and a spuriously correlated background is a usable shortcut
for doing so. SAM\,3 is told what to find. The prompt \texttt{"impala"} supplies the foreground identity and
removes the incentive to ride the scene --- exactly the direction \citet{xiao}'s own finding predicts, since
better and less shortcut-prone models close the background gap. PanAf's target compounds the difference, since
fine-grained behaviour is intrinsically background-correlated whereas finding and following a named animal is
not.

The result is therefore an extension rather than a refutation. The backgrounds that causally drive a
camera-trap classifier's location drop do not, once isolated from the animal and used as a before-running
distance, forecast a promptable tracker's transfer. Because we isolate scene from animal explicitly, a step
none of the label-free-prediction methods take, this is a cleaner statement than ``background is diagnostic''
rather than a weaker one. The claim is bounded accordingly: it concerns background distance as a predictor of
detection transfer for a promptable tracker, not a blanket claim that background is irrelevant.

\section{What this means in practice}
\label{sec:disc-practice}
The field problem that motivated this work --- telling a practitioner, before spending compute, whether a
promptable tracker is trustworthy for a new species and place --- is not solved by label-free distances. That
is a negative answer, but a useful one. The intuition that far-from-training implies unreliable does not hold
here, so a deployment gate cannot be built from taxonomic, visual, environmental or temporal distance to a
reference set. The one predictive label-free signal, animal size, is too coarse and too obvious to serve as
one. What the negative result does establish is a baseline: any future before-running estimator must clear the
same whole-group hold-out bar that these well-motivated distances failed.

One after-running signal does clear that bar, and it is stated plainly here rather than left implicit. If a
practitioner is willing to run the tracker first, an ATC-style estimator~\cite{atc} --- reading SAM\,3's own
detection confidence on the target and thresholding it --- predicts detection out of sample on both
cross-validation schemes, where every label-free distance failed. It is deliberately kept outside the
contribution, for three reasons. It is after-running, a weaker operating point than the before-running question
this work set out to answer. It is confirmatory of an established family rather than novel. And it is
mechanistically close to circular, since feature and target are both functions of SAM\,3's own output: a cell
with no confident masklet trivially has zero detection accuracy. It is a real, validated estimator for the
after-running, before-labelling operating point --- but not the before-running predictor this dissertation asks
for, and not part of the null it reports.

\section{Limitations and threats to validity}
\label{sec:limitations}

\paragraph{The reference is not the pretraining corpus.} Distances are measured against the SA-FARI train
split, not against SAM\,3's true pretraining data, which is undisclosed. The null is scoped accordingly. One
hedge was tested: the familiarity proxy reads separability directly from SAM\,3's own features
(Section~\ref{sec:familiarity}), and it too is null, recovering the same size-entangled visual signal. Even
model-internal familiarity does not lift the result. A coverage-anchored proxy for the actual corpus remains
untested.

\paragraph{One dataset and one backbone.} A single inference pass is reused across both splits, so the species
and location experiments are two overlays on the same scored cells rather than independent replications. Two
steps were taken against this. Against the dataset threat, the whole pipeline was re-run on an independent
dataset, where the null holds again (Section~\ref{sec:burst}) --- a second convergent null rather than positive
proof, since that dataset remains underpowered for a small species-constant effect. Against the model threat,
the controlled swap put two architecturally distinct trackers through the same procedure with the same result
(Section~\ref{sec:model_swap}).

A third dataset, MammAlps~\cite{mammalps}, was run in an earlier pass. It is directionally consistent with the null but is not
counted here: at roughly $20$ cells over $5$ species it is too small to certify one, and its per-clip
annotations are not reproducible from the public release. It is reported as excluded rather than silently
dropped.

\paragraph{The effective sample is groups, not cells.} Tens of independent species and locations, not hundreds
of cells. That is why inference is group-clustered throughout, and why aggregate correlations are reported
alongside their much weaker cell-level values, to avoid the ecological-correlation fallacy.

\paragraph{Representational probing is only partly opened.} The familiarity proxy tests whether SAM\,3's
features separate the species. They do, but that separability merely re-encodes the visual distance and does
not predict transfer. A fuller probe of what those representations encode remains open.

\paragraph{The scorer.} All scores come from the unmodified official evaluator, which wraps the standard
HOTA-family metrics rather than any SAM-3-specific logic. Its one SAM-3-calibrated knob is a fixed $0.5$
confidence gate, and Section~\ref{sec:model_swap} establishes that this gate is not load-bearing for the
cross-model comparison. The measurement is therefore model-fair.

\section{Ethical considerations}
\label{sec:disc-ethics}
The data are anonymised camera-trap footage released under a non-commercial licence, and location identifiers
are already anonymised in SA-FARI. The negative result carries a responsible-deployment message. Because
label-free distance does not reliably flag when a promptable tracker will fail, conservation teams should not
treat low distance from training data as a safety guarantee. Reliability should be checked with labelled
spot-audits rather than assumed.
